\section{Balance-of-payments booking of the adoption flows}
\label{app:booking}

The sign and magnitude of the adoption channel's contribution to output depend
on how the green-margin resource flows---the switching cost $\psi_g D^{sw}$ and
the green energy bill---are booked in the balance of payments. The main text
reports the import booking throughout. This appendix collects the two results
that hinge on that choice: the domestic-booking sign comparison, and the
carbon-revenue recycling experiment, which is naturally a booking-conditional
statement about where the switching subsidy lands.

\subsection{Import versus domestic booking of the adoption flows}\label{app:booking_flows}
Both statics of Section~\ref{sec:signmap} are conditional on the import booking of the
adoption flows. We solve the model under the alternative domestic booking of
Section~\ref{sec:bop}---green energy and green installation as competitive,
zero-profit domestic value added rebated to households, with only brown energy
imported---which resolves the phantom-dividend problem and clears to machine
precision. Under this booking neither the switching cost nor the green energy
bill is a leakage, and the adoption channel's contribution to cumulative output
($\Delta_{24}$) is as following:
\IfFileExists{output/tab_signmap_booking.tex}{\begin{table}[H]\centering
\small
\begin{tabular}{lcc}
\toprule
$\Delta_{24}$, adoption output contribution & import booking & domestic booking\\
\midrule
price shock & $-1.60$ & $+35.98$\\
supply shock & $+0.38$ & $+5.96$\\
\bottomrule
\end{tabular}
\caption{Adoption channel contribution to cumulative output ($\Delta_{24}$), import vs.\ domestic booking (Section~\ref{sec:bop}). \emph{None} policy.}
\label{tab:signmap}
\end{table}
}


\subsection{Recycling the carbon revenue: rebate versus green subsidy}
\label{sec:recycle}

The welfare-optimal carbon tax of Section~\ref{sec:carbontax} returns its
revenue lump sum. An alternative earmarks it to the transition itself, paying a
share $s_g$ of the green switching cost. We compare the two recyclings across
carbon-tax rates $\tau_b\in\{0.10,0.20,0.30\}$, holding $\psi_g$ at its no-tax
value so the steady-state green share floats (Figures~\ref{fig:taub_rebate}
and~\ref{fig:taub_green}). Under the lump-sum rebate the steady-state green
share rises from $8.9\%$ to $20.1\%$ across the grid, and the whole brown-price
path in the crisis shifts up by the factor $(1+\tau_b)$: the pre-tax price
passes the world shock through identically, while households pay the
tax-inclusive price. Under \emph{budget-neutral} green-subsidy recycling---$s_g$
solved so that the switching subsidy exactly exhausts the carbon revenue
($T^{\mathrm{reb}}=0$)---the revenue funds a subsidy covering $35$--$40\%$ of
the switching cost and lifts the steady-state green share to $27$--$46\%$.

Two features are worth flagging. First, the carbon base is self-eroding: at
$\tau_b=0.10$ the revenue is a third smaller under green-subsidy recycling than
under the rebate, because the greener economy holds less brown-durable stock to
tax. Second, the crisis-time greening response is \emph{smaller}, not larger,
from the greener starting point under the subsidy: recycling that greens the
economy through a lower switching \emph{cost} preferentially converts the
households nearest the logit margin, leaving a residual brown pool that is less
responsive to the shock---the mirror image of the rebate, which greens through a
higher permanent brown \emph{price} and leaves the residual pool closer to the
margin. This is the adoption-absorber exhaustion of Section~\ref{sec:routeA}
seen across policy instruments rather than across shock persistence.

\begin{figure}[H]\centering
\includegraphics[width=\textwidth]{output/fig_exante_taub_rebate.png}
\caption{Carbon-tax sweep, lump-sum rebate: laissez-faire and ex-post cap
against the ex-ante ETS economy at $\tau_b\in\{0.10,0.20,0.30\}$.}
\label{fig:taub_rebate}
\end{figure}

\begin{figure}[H]\centering
\includegraphics[width=\textwidth]{output/fig_exante_taub_greensub_bn.png}
\caption{Carbon-tax sweep, budget-neutral green-subsidy recycling
($T^{\mathrm{reb}}=0$, $s_g$ endogenous). Same axes as
Figure~\ref{fig:taub_rebate}.}
\label{fig:taub_green}
\end{figure}
